\subsection{Orthogonality, Radicals, and Nonsingular Metric Spaces}

Throughout this subsection, metric vector spaces are finite-dimensional over an arbitrary field $\mathbb{F}$.

\begin{definition} \label{def:metric-orthogonality-isotropy-radical} \hlt{Orthogonality, Isotropy, and Radical}\\
Let $\mathcal{V}$ be a metric vector space. By Definition \ref{def:metric-vector-space-basic-form-types}(ii)--(iv) and Proposition \ref{prop:form-type-relations-by-characteristic}(i), vanishing of the form is symmetric in its two arguments, so left and right orthogonality agree.
\begin{enumerate}[label=(\roman*)]
\setlength{\itemsep}{0pt}
\item $\mathbf{x}\perp\mathbf{y}$ means $\innerprod{\mathbf{x},\mathbf{y}}=0$; then $\mathbf{x},\mathbf{y}$ are \hlt{orthogonal}. For $S\subseteq\mathcal{V}$, $S^{\perp}\defeq\{\mathbf{v}\in\mathcal{V}:(\forall\,\mathbf{s}\in S)\,\innerprod{\mathbf{v},\mathbf{s}}=0\}$.
\item $\mathcal{S}\leq\mathcal{V}$ is \hlt{totally isotropic} or \hlt{totally degenerate} if $\mathcal{S}\subseteq\mathcal{S}^{\perp}$, and \hlt{nonsingular as a subspace} if $\mathcal{S}\cap\mathcal{S}^{\perp}=\{\mathbf{0}\}$.
\item $\mathbf{x}\neq\mathbf{0}$ is \hlt{isotropic} or \hlt{null} if $\innerprod{\mathbf{x},\mathbf{x}}=0$, and \hlt{nonisotropic} otherwise.\\
The space $\mathcal{V}$ is \hlt{isotropic} if $(\exists\,\mathbf{x}\in\mathcal{V}\setminus\{\mathbf{0}\})\,\innerprod{\mathbf{x},\mathbf{x}}=0$, and \hlt{anisotropic} if $(\forall\,\mathbf{x}\in\mathcal{V}\setminus\{\mathbf{0}\})\,\innerprod{\mathbf{x},\mathbf{x}}\neq0$.
\item $\mathbf{x}$ is \hlt{degenerate} if $(\forall\,\mathbf{v}\in\mathcal{V})\,\innerprod{\mathbf{x},\mathbf{v}}=0$. The \hlt{radical} is $\operatorname{rad}(\mathcal{V})\defeq\mathcal{V}^{\perp}$.
\item $\mathcal{V}$ is \hlt{nonsingular} or \hlt{nondegenerate} if $\operatorname{rad}(\mathcal{V})=\{\mathbf{0}\}$, \hlt{singular} otherwise.
\item $\mathcal{V}$ is \hlt{totally singular} or \hlt{totally degenerate} if $\operatorname{rad}(\mathcal{V})=\mathcal{V}$.
\end{enumerate}
\end{definition}

\begin{proposition} \label{prop:metric-radical-equals-bilinear-radicals} \hlt{Metric Radical as a Bilinear Radical}\\
Let $B(\mathbf{x},\mathbf{y})=\innerprod{\mathbf{x},\mathbf{y}}$ be the form of a metric vector space $\mathcal{V}$.
\begin{enumerate}[label=(\roman*)]
\setlength{\itemsep}{0pt}
\item $\operatorname{rad}_{L}(B)=\operatorname{rad}_{R}(B)=\operatorname{rad}(\mathcal{V})$.
\item $\mathcal{V}$ is nonsingular $\Longleftrightarrow$ $B$ is perfect.
\end{enumerate}
\begin{proof}
\begin{enumerate}[label=(\roman*)]
\setlength{\itemsep}{0pt}
\item[$(i)$] By the definitions of $\mathcal{V}^{\perp}$ and $\operatorname{rad}_{L}(B)$, $\operatorname{rad}(\mathcal{V})=\operatorname{rad}_{L}(B)$. The equality $\operatorname{rad}_{L}(B)=\operatorname{rad}_{R}(B)$ holds by Proposition \ref{prop:form-type-relations-by-characteristic}(iv).
\item[$(ii)$] ($\Rightarrow$) If $\mathcal{V}$ is nonsingular, then $\operatorname{rad}(\mathcal{V})=\{\mathbf{0}\}$, so $\operatorname{rad}_{L}(B)=\{\mathbf{0}\}$ by (i). Hence $B$ is perfect by Proposition \ref{prop:square-bilinear-pairing-equivalences}.
\item[$(ii)$] ($\Leftarrow$) If $B$ is perfect, then $\operatorname{rad}_{L}(B)=\{\mathbf{0}\}$ by Proposition \ref{prop:square-bilinear-pairing-equivalences}. Hence $\operatorname{rad}(\mathcal{V})=\{\mathbf{0}\}$ by (i), so $\mathcal{V}$ is nonsingular.
\end{enumerate}
\hlt{Structural Commentary}: Radical identification reduces perfection to triviality of the common radical.\\
\hlt{Key Insight}: A metric form is perfect exactly when its common radical vanishes.
\end{proof}
\end{proposition}

Let $\beta$ be an ordered basis of $\mathcal{V}$. By Proposition \ref{prop:metric-radical-equals-bilinear-radicals}(i) and Proposition \ref{prop:bilinear-radicals-matrix-and-dual-maps}, the coordinate radical is $\ker(M_{\beta}^{t})$ and $\dim\operatorname{rad}(\mathcal{V})=\dim\mathcal{V}-\operatorname{rank}(M_{\beta})$. By Proposition \ref{prop:metric-radical-equals-bilinear-radicals}(ii) and Proposition \ref{prop:square-bilinear-pairing-equivalences}, $\mathcal{V}$ is nonsingular $\Longleftrightarrow$ $M_{\beta}$ is invertible.

\begin{remark} \label{rem:subspace-radical-versus-ambient-complement} \hlt{Subspace Radicals}\\
For $\mathcal{S}\leq\mathcal{V}$, the restricted form satisfies $\operatorname{rad}(\mathcal{S})=\{\mathbf{s}\in\mathcal{S}:(\forall\,\mathbf{t}\in\mathcal{S})\,\innerprod{\mathbf{s},\mathbf{t}}=0\}=\mathcal{S}\cap\mathcal{S}^{\perp}$.\\
In general, $\operatorname{rad}(\mathcal{S})\neq\mathcal{S}^{\perp}$ because the latter is computed in the ambient space.
\end{remark}

\begin{proposition} \label{prop:riesz-map-isomorphism-criterion} \hlt{Riesz Map Isomorphism Criterion}\\
Let $\mathcal{V}$ be a metric vector space and define $\tau:\mathcal{V}\to\mathcal{V}'$ by $\tau(\mathbf{x})(\mathbf{v})=\innerprod{\mathbf{v},\mathbf{x}}$.
\begin{enumerate}[label=(\roman*)]
\setlength{\itemsep}{0pt}
\item $\tau\in\LL(\mathcal{V},\mathcal{V}')$ and $\ker(\tau)=\operatorname{rad}(\mathcal{V})$.
\item $\tau$ is an isomorphism $\Longleftrightarrow$ $\operatorname{rad}(\mathcal{V})=\{\mathbf{0}\}$.
\end{enumerate}
\begin{proof}
\begin{enumerate}[label=(\roman*)]
\setlength{\itemsep}{0pt}
\item[$(i)$] For the metric form $B$, $\tau=\Psi_B$, so $\tau\in\LL(\mathcal{V},\mathcal{V}')$ by Proposition \ref{prop:bilinear-forms-linearize-through-duals}(i).\\
The identity $\ker(\tau)=\operatorname{rad}_{R}(B)=\operatorname{rad}(\mathcal{V})$ follows by Proposition \ref{prop:bilinear-radicals-matrix-and-dual-maps}(i) and Proposition \ref{prop:metric-radical-equals-bilinear-radicals}(i).
\item[$(ii)$] ($\Rightarrow$) If $\tau$ is isomorphism, then it is injective by Proposition \ref{prop:linear-map-invertible-iff-bijective}, so $\ker(\tau)=\{\mathbf{0}\}$ by Proposition \ref{prop:injective-iff-kernel-zero}; $\operatorname{rad}(\mathcal{V})=\{\mathbf{0}\}$ by (i).
\item[$(ii)$] ($\Leftarrow$) If $\operatorname{rad}(\mathcal{V})=\{\mathbf{0}\}$, then $\ker(\tau)=\{\mathbf{0}\}$ by (i), so $\tau$ is injective by Proposition \ref{prop:injective-iff-kernel-zero}. By Theorem \ref{thm:dual-basis-and-dimension}(ii), $\dim(\mathcal{V})=\dim(\mathcal{V}')$, so $\tau$ is surjective by Corollary \ref{cor:equal-dimension-kernel-range}(ii) and is an isomorphism by Proposition \ref{prop:linear-map-invertible-iff-bijective}.
\end{enumerate}
\hlt{Structural Commentary}: Kernel identification and equal dimensions prove the isomorphism criterion.\\
\hlt{Key Insight}: The Riesz map is invertible exactly when the metric form is nonsingular.
\end{proof}
\end{proposition}

\begin{theorem} \label{thm:riesz-map-for-nonsingular-metric-spaces} \hlt{Riesz Map for Nonsingular Metric Spaces}\\
Let $\mathcal{V}$ be a nonsingular metric vector space. Then $(\forall\,f\in\mathcal{V}')(\exists!\,\mathbf{x}\in\mathcal{V})(\forall\,\mathbf{v}\in\mathcal{V})\,f(\mathbf{v})=\innerprod{\mathbf{v},\mathbf{x}}$.
\begin{proof}
Let $\tau:\mathcal{V}\to\mathcal{V}'$ be defined by $\tau(\mathbf{x})(\mathbf{v})=\innerprod{\mathbf{v},\mathbf{x}}$. By Proposition \ref{prop:riesz-map-isomorphism-criterion}(ii), $\tau$ is an isomorphism.\\
Its bijectivity is precisely the stated existence and uniqueness.\\
\hlt{Structural Commentary}: The Riesz isomorphism turns representation into unique inversion.\\
\hlt{Key Insight}: Nonsingularity, rather than positivity, is the exact hypothesis for algebraic Riesz representation.
\end{proof}
\end{theorem}

\begin{definition} \label{def:metric-isometry-and-classical-groups} \hlt{Isometries and Classical Groups}\\
Let $\mathcal{V},\mathcal{W}$ be metric vector spaces.
\begin{enumerate}[label=(\roman*)]
\setlength{\itemsep}{0pt}
\item An \hlt{isometry} $\sigma:\mathcal{V}\to\mathcal{W}$ is a bijective linear map satisfying $(\forall\,\mathbf{x},\mathbf{y}\in\mathcal{V})\,\innerprod{\sigma\mathbf{x},\sigma\mathbf{y}}=\innerprod{\mathbf{x},\mathbf{y}}$; write $\mathcal{V}\cong\mathcal{W}$ if one exists.
\item The group of isometries $\mathcal{V}\to\mathcal{V}$ is denoted $\operatorname{Iso}(\mathcal{V})$. If $\mathcal{V}$ is nonsingular orthogonal, this is the \hlt{orthogonal group} $\mathcal{O}(\mathcal{V})$.\\
If $\mathcal{V}$ is nonsingular symplectic, this is the \hlt{symplectic group} $\operatorname{Sp}(\mathcal{V})$.
\end{enumerate}
\end{definition}

\begin{proposition} \label{prop:matrix-criterion-for-metric-isometries} \hlt{Matrix Criterion for Isometries}\\
Let $\beta$ be an ordered basis of a metric vector space $\mathcal{V}$, let $\sigma\in\LL(\mathcal{V})$, let $M=M_{\beta}$, and let $A=[\sigma]_{\beta\leftarrow\beta}$.
\begin{enumerate}[label=(\roman*)]
\setlength{\itemsep}{0pt}
\item $\sigma\in\operatorname{Iso}(\mathcal{V})$ $\Longleftrightarrow$ $\sigma$ is bijective and $A^tMA=M$.
\item If $\sigma\in\operatorname{Iso}(\mathcal{V})$ and $M$ is invertible, then $\det A=\pm1$.
\end{enumerate}
\begin{proof}
\begin{enumerate}[label=(\roman*)]
\setlength{\itemsep}{0pt}
\item[$(i)$] ($\Rightarrow$) Suppose that $\sigma\in\operatorname{Iso}(\mathcal{V})$. Then $\sigma$ is bijective.\\
By Proposition \ref{prop:metric-form-matrix-recovery-and-congruence}(i), metric preservation implies $(\forall\,X,Y\in\mathbb{F}^{n})\,X^tA^tMAY=X^tMY$, where $n=\dim\mathcal{V}$.\\
Taking $X,Y$ to be standard coordinate columns proves $A^tMA=M$.
\item[$(i)$] ($\Leftarrow$) Suppose that $\sigma$ is bijective and $A^tMA=M$. By Proposition \ref{prop:metric-form-matrix-recovery-and-congruence}(i), $(\forall\,\mathbf{x},\mathbf{y}\in\mathcal{V})\,\innerprod{\sigma\mathbf{x},\sigma\mathbf{y}}=\innerprod{\mathbf{x},\mathbf{y}}$. Thus $\sigma\in\operatorname{Iso}(\mathcal{V})$.
\item[$(ii)$] Suppose that $\sigma\in\operatorname{Iso}(\mathcal{V})$ and $M$ is invertible. By (i), $A^tMA=M$. Hence $\det(M)=\det(A^tMA)=(\det A)^2\det(M)$.\\
Since $\det(M)\neq0$, $(\det A)^2=1$, so $\det A=\pm1$.
\end{enumerate}
\hlt{Structural Commentary}: The matrix identity and determinants prove both claims.\\
\hlt{Key Insight}: The equation $A^tMA=M$ realizes an isometry group as a congruence stabilizer.
\end{proof}
\end{proposition}

\begin{definition} \label{def:metric-adjoint-nonsingular-form} \hlt{Adjoint for a Nonsingular Bilinear Form}\\
Let $\mathcal{V}$ be nonsingular. For $T\in\LL(\mathcal{V})$, an operator $S\in\LL(\mathcal{V})$ is a \hlt{metric adjoint} of $T$ if $(\forall\,\mathbf{x},\mathbf{y}\in\mathcal{V})\,\innerprod{T\mathbf{x},\mathbf{y}}=\innerprod{\mathbf{x},S\mathbf{y}}$.\\
When this operator is unique, it is denoted by $T^*$.
\end{definition}

\begin{proposition} \label{prop:metric-adjoints-and-isometries} \hlt{Metric Adjoints and Isometries}\\
Let $\mathcal{V}$ be nonsingular, let $B(\mathbf{x},\mathbf{y})=\innerprod{\mathbf{x},\mathbf{y}}$ be its form, and let $T\in\LL(\mathcal{V})$.
\begin{enumerate}[label=(\roman*)]
\setlength{\itemsep}{0pt}
\item $(\exists!\,T^*\in\LL(\mathcal{V}))(\forall\,\mathbf{x},\mathbf{y}\in\mathcal{V})\,\innerprod{T\mathbf{x},\mathbf{y}}=\innerprod{\mathbf{x},T^*\mathbf{y}}$.
\item If $\beta$ is an ordered basis, $M=M_{\beta}$, and $A=[T]_{\beta}$, then $[T^*]_{\beta}=M^{-1}A^{t}M$.
\item $T$ is an isometry $\Longleftrightarrow$ $T$ is bijective and $T^*T=I$.
\item If $T$ is bijective, then $T^*T=I$ $\Longleftrightarrow$ $T^*=T^{-1}$.
\end{enumerate}
\begin{proof}
\begin{enumerate}[label=(\roman*)]
\setlength{\itemsep}{0pt}
\item[$(i)$] Define $\tau:\mathcal{V}\to\mathcal{V}'$ by $\tau(\mathbf{y})(\mathbf{x})=\innerprod{\mathbf{x},\mathbf{y}}$ and $F:\mathcal{V}\to\mathcal{V}'$ by $F(\mathbf{y})(\mathbf{x})=\innerprod{T\mathbf{x},\mathbf{y}}$.\\
By Proposition \ref{prop:riesz-map-isomorphism-criterion}(ii), $\tau$ is an isomorphism, while bilinearity implies $F\in\LL(\mathcal{V},\mathcal{V}')$. Set $T^*=\tau^{-1}\circ F\in\LL(\mathcal{V})$.\\
Then $(\forall\,\mathbf{x},\mathbf{y}\in\mathcal{V})\,\innerprod{T\mathbf{x},\mathbf{y}}=\innerprod{\mathbf{x},T^*\mathbf{y}}$. If $S\in\LL(\mathcal{V})$ satisfies this identity, then $\tau\circ S=F=\tau\circ T^*$, so $S=T^*$.
\item[$(ii)$] By Proposition \ref{prop:metric-radical-equals-bilinear-radicals}(ii) and Proposition \ref{prop:square-bilinear-pairing-equivalences}, $M$ is invertible. Put $D=[T^*]_{\beta}$.\\
By the adjoint identity and Proposition \ref{prop:metric-form-matrix-recovery-and-congruence}(i), $(\forall\,X,Y\in\mathbb{F}^{n})\,X^tA^tMY=X^tMDY$.\\
Taking $X,Y$ to be standard coordinate columns proves $A^tM=MD$, so $D=M^{-1}A^tM$.
\item[$(iii)$] ($\Rightarrow$) Suppose that $T$ is an isometry. Then $T$ is bijective. Fix $\mathbf{x},\mathbf{y}\in\mathcal{V}$.\\
The adjoint identity and metric preservation imply $\innerprod{\mathbf{x},T^*T\mathbf{y}}=\innerprod{T\mathbf{x},T\mathbf{y}}=\innerprod{\mathbf{x},\mathbf{y}}$.\\
Hence $(T^*T-I)\mathbf{y}\in\operatorname{rad}_{R}(B)=\{\mathbf{0}\}$ by Proposition \ref{prop:metric-radical-equals-bilinear-radicals}(i), and therefore $T^*T=I$.
\item[$(iii)$] ($\Leftarrow$) Suppose that $T$ is bijective and $T^*T=I$. For $\mathbf{x},\mathbf{y}\in\mathcal{V}$, the adjoint identity implies $\innerprod{T\mathbf{x},T\mathbf{y}}=\innerprod{\mathbf{x},T^*T\mathbf{y}}=\innerprod{\mathbf{x},\mathbf{y}}$. QED.
\item[$(iv)$] ($\Rightarrow$) If $T^*T=I$, then right-composition with $T^{-1}$ proves $T^*=T^{-1}$.
\item[$(iv)$] ($\Leftarrow$) If $T^*=T^{-1}$, then $T^*T=T^{-1}T=I$.
\end{enumerate}
\hlt{Structural Commentary}: Riesz duality constructs the adjoint; coordinate and inverse identities finish.\\
\hlt{Key Insight}: Nonsingularity makes the adjoint unique, and a bijective operator is an isometry exactly when its adjoint is its inverse.
\end{proof}
\end{proposition}

\begin{proposition} \label{prop:metric-spaces-isometric-iff-congruent-matrices} \hlt{Isometry Is Matrix Congruence}\\
Let $(\mathcal{V},B_{\mathcal{V}})$ and $(\mathcal{W},B_{\mathcal{W}})$ be $n$-dimensional metric vector spaces with ordered bases $\beta=(\mathbf{v}_1,\ldots,\mathbf{v}_n)$ and $\gamma=(\mathbf{w}_1,\ldots,\mathbf{w}_n)$ and representing matrices $M_{\beta}$ and $N_{\gamma}$. $\mathcal{V}\cong\mathcal{W}$ $\Longleftrightarrow$ $M_{\beta}$ and $N_{\gamma}$ are congruent.
\begin{proof}
\begin{enumerate}[label=\roman*.]
\setlength{\itemsep}{0pt}
\item[$(\Rightarrow)$] Let $\sigma:\mathcal{V}\to\mathcal{W}$ be an isometry. It is a bijective linear map and hence an isomorphism by Proposition \ref{prop:linear-map-invertible-iff-bijective}.\\
Thus $\delta=(\sigma^{-1}\mathbf{w}_1,\ldots,\sigma^{-1}\mathbf{w}_n)$ is a basis of $\mathcal{V}$ by Proposition \ref{prop:isomorphisms-preserve-bases}.\\
Metric preservation implies $(\forall\,i,j\in\{1,\ldots,n\})\,(M_{\delta})_{ij}=(N_{\gamma})_{ij}$, so $M_{\delta}=N_{\gamma}$.\\
Hence $N_{\gamma}$ is congruent to $M_{\beta}$ by Proposition \ref{prop:metric-form-matrix-recovery-and-congruence}(iii).
\item[$(\Leftarrow)$] Suppose that $M_{\beta}$ and $N_{\gamma}$ are congruent. By Proposition \ref{prop:metric-form-matrix-recovery-and-congruence}(iii), fix an ordered basis $\delta=(\mathbf{d}_1,\ldots,\mathbf{d}_n)$ of $\mathcal{V}$ satisfying $M_{\delta}=N_{\gamma}$.\\
By Proposition \ref{prop:linear-map-determined-by-basis}, fix linear maps $\sigma:\mathcal{V}\to\mathcal{W}$ and $\rho:\mathcal{W}\to\mathcal{V}$ satisfying $(\forall\,j\in\{1,\ldots,n\})\,\sigma\mathbf{d}_j=\mathbf{w}_j$ and $\rho\mathbf{w}_j=\mathbf{d}_j$.\\
The compositions $\rho\sigma$ and $\sigma\rho$ agree with respective identity maps on $\delta$ and $\gamma$, so $\rho\sigma=I_{\mathcal{V}}$ and $\sigma\rho=I_{\mathcal{W}}$ by Proposition \ref{prop:linear-map-determined-by-basis}.\\
Hence $\sigma$ is invertible with inverse $\rho$ and is bijective by Proposition \ref{prop:linear-map-invertible-iff-bijective}. Fix $\mathbf{x},\mathbf{y}\in\mathcal{V}$ with $\delta$-coordinate columns $X,Y$.\\
These are the $\gamma$-coordinate columns of $\sigma\mathbf{x},\sigma\mathbf{y}$. By Proposition \ref{prop:metric-form-matrix-recovery-and-congruence}(i), the form values are $X^tN_{\gamma}Y$ and $X^tM_{\delta}Y$, respectively.\\
They agree because $M_{\delta}=N_{\gamma}$, so $\sigma$ is an isometry.
\end{enumerate}
\hlt{Structural Commentary}: Basis transport proves both directions of the congruence criterion.\\
\hlt{Key Insight}: Isometry classes of metric spaces coincide with congruence classes of their representing matrices.
\end{proof}
\end{proposition}

\begin{perspective} \label{per:ordinary-congruence-bilinear-form-bridge} \hlt{Congruence versus Similarity}\\
Similarity $N=P^{-1}MP$ governs operator matrices. Congruence $N=P^tMP$ governs bilinear-form matrices because the same basis change acts in both slots. Over $\mathbb{R}$, transpose equals conjugate transpose, so real form classification uses congruence.
\end{perspective}
